\documentclass[a4paper]{article}

%% Language and font encodings
\usepackage[english]{babel}
\usepackage[utf8x]{inputenc}
\usepackage[T1]{fontenc}
\usepackage{blindtext}
\usepackage{enumitem}
\usepackage{float}
\usepackage{commath}

%% Sets page size and margins
\usepackage[a4paper,top=3cm,bottom=2cm,left=3cm,right=3cm,marginparwidth=1.75cm]{geometry}

%% Useful packages
\usepackage{amsmath}
\usepackage{graphicx}
\usepackage[colorinlistoftodos]{todonotes}
\usepackage[colorlinks=true, allcolors=blue]{hyperref}

\title{Práctica Laboratorio de Aerodinámica No. XX \\
Plantilla de informe de laboratorio}

\author{Juan Carlos Perafan, Juan Pablo Alvarado, Daniel Restrepo}

\date{\today}

\begin{document}
\maketitle

\section{Introducción}

La elaboración del informe de práctica de laboratorio tiene como finalidad desarrollar en el alumno la habilidad de expresarse de manera escrita en forma clara y concisa. En dicho informe se incluirán: los fundamentos teóricos; el procedimiento y las características de los métodos y materiales empleados; la observación y manipulación de datos para la obtención de resultados; la presentación, interpretación y análisis de dichos resultados, así como la relación que tiene con la teoría; sus implicaciones, limitaciones de las experiencias y posibles sugerencias.

Los informes han de estar estructurados de la siguiente forma:

\begin{itemize}
\item \textbf{Práctica Laboradorio de Aerodinámica No. XX}
\item \textbf{Título}
\item \textbf{Nombres/ID}
\item \textbf{Fecha}
\item \textbf{Preinforme – $20\%$}
\item \textbf{Metodología – $10\%$}
\item \textbf{Análisis de resultados – $30\%$}
\item \textbf{Conclusiones – $30\%$}
\item \textbf{Referencias Bibliográficas – $5\%$}
\item \textbf{Anexos (En caso de tenerlos)}
\end{itemize}

\textit{*Nota 1:} con el 5\% restante se califica la redacción y la ortografía del informe.

\textit{*Nota 2:} El docente encargado determinará el método de entrega del informe y preinforme (físico, correo electrónico, a través de Overleaf, etc.) y las especificaciones de esta (formato del archivo, presentación, etc.)

Se resalta que los informes de laboratorio no tienen hoja de presentación, es decir en la primera hoja se deben consignar los primeros 4 ítems de la lista anterior para luego comenzar en la misma hoja con el resto (Preinforme, Metodología, etc.). A continuación, se analizarán cada una de estas partes por separado enfatizando los aspectos más importantes de cada uno de ellos.

\section{Preinforme}
El preinforme de cada práctica consta de dos partes: la planeación y los cálculos previos.

La primera parte consiste en explicar de manera breve la planeación previa para la realización de la práctica, indicando todo lo que se tuvo en cuenta para su ejecución, (e.g.: materiales, tiempos, modelos, tablas). Se busca que el estudiante adquiera la habilidad de poder diseñar, estructurar y programar una sesión experimental, teniendo en cuenta los diferentes aspectos que esto conlleva.

La segunda parte consiste en realizar los cálculos que se exigen en las guías de laboratorio, estos suelen ser usados como referencia teórica para los cálculos del error en los experimentos. Además de esto, en esta parte también se debe reportar los cálculos previos sobre los modelos a ensayar, (e.g.: calcular la sustentación y el arrastre del ala finita diseñada), y así determinar qué balanza se debe usar en el experimento.

\section{Procedimiento (Datos experimentales)}
Explicación más profunda del procedimiento realizado en la práctica de laboratorio, \textbf{indicando el paso a paso detallado: qué, cómo, y para qué se tomaron los datos.} Se deben mostrar todas las medidas directas que se realizan durante la práctica para posterior análisis.

Se debe reportar el manejo y tratamiento que se le dio a los datos y observaciones experimentales para llegar a los resultados. El detalle con el cual se deben presentar los procedimientos de manipulación de los datos obtenidos en cada práctica dependerá de los objetivos específicos de las mismas y se indicará en cada una de las guías.

Los resultados deben presentarse de manera clara, por lo general en forma de figuras o tablas y con su respectivo grado de validez.

 Resulta importante establecer qué tipo de error se comete en ellas y como se propaga durante su manipulación para la obtención de resultados. Los datos experimentales deben ser presentados en forma de tablas o de gráficas dependiendo de la forma de salida de los datos experimentales.


\section{Análisis de Resultados}
Esta sección es una de las más importantes del informe. Aquí se debe presentar los resultados obtenidos de manera clara y concisa.

Resulta importante establecer qué tipo de error se comete en la toma de datos y cómo este se propaga durante su manipulación para la obtención de resultados. Los datos experimentales deben ser presentados en forma de tablas o de gráficas dependiendo de la forma de salida de los datos experimentales.

Se debe explicar el comportamiento de los resultados obtenidos y lo que reflejan al compararlos con los datos teóricos (presentados en el preinforme). El alumno debe desplegar su capacidad de análisis, relacionando causas y efectos, comparando unos resultados con los otros obtenidos en la experiencia actual, experiencias anteriores y teoría. Los errores propios del trabajo experimental es correcto comentarlos, pero no deben transformarse en el centro del análisis.

\section{Conclusiones}
Deben darse de forma clara y técnica y solo se deben referir a los procedimientos, resultados y análisis elaborados en la práctica. Conclusiones con respecto a los resultados obtenidos en la práctica de laboratorio y análisis de los errores encontrados. 

En esta sección se incluyen los conocimientos adquiridos durante la experiencia práctica, considerando los conocimientos teóricos que se poseían previamente y los aspectos aprendidos en la sección práctica. Se emitirá opinión acerca de la utilidad, la viabilidad, aplicabilidad, ventajas, así como sobre las limitaciones de la técnica o del procedimiento realizado. Si las hubiere, recomendaciones para la optimización de la experiencia práctica. Es importante resaltar que la extensión de las conclusiones esta limitada por las características del procedimiento experimental y las variables evaluadas durante el experimento. No deben incluirse conclusiones que involucren variables no evaluadas o que sobrepasen el marco de la experiencia realizada.

En general cada una de las conclusiones debe estar relacionada ya sea con los objetivos generales o con los objetivos específicos. Estas son las primeras que deben ser enunciadas. Enseguida, deben enunciarse aquellas conclusiones que no estén directamente relacionadas con los objetivos generales y/o específicos. \textbf{Debe cuidarse de no confundir las conclusiones con los resultados.}

\section{Referencias}
En esta parte se incluyen todas las fuentes bibliográficas consultadas para la elaboración del informe. Las referencias deben numerarse de tal modo que se determine de una forma fácil en que parte del texto se está utilizando mediante el método de referencia cruzada. Las notas del curso o lo dicho en las clases o la información encontrada en las presentaciones y guías de laboratorio no se pueden tomar como una referencia para estos informes.
Cualquier idea o concepto que no sea de su propiedad, debe ser referenciado, no son aceptadas referencias de Wikipedia, Rincón del Vago, etc. sean consistentes con el estilo de referencia. A continuación se muestra un ejemplo de referencia cruzada de acuerdo a la información suministrada en el archivo sample.bib. Para llamar una referencia se debe utilizar el comando \textit{cite} y a continuación el nombre de la referencia en este caso \textit{greenwade93}, \cite{greenwade93}, el título References, aparecerá automáticamente cuando se utilizan los comandos requeridos en esta sección: 

 \bibliographystyle{IEEEtran}
\bibliography{sample}

\section{Anexos}
Aquí se incluye toda aquella información auxiliar o complementaria para los lectores. En esta sección se incluyen los cálculos requeridos para la obtención de los resultados, los ejemplos de cálculo para cada una de las variables presentadas en los resultados, así como de las variables involucradas en los pasos intermedios.

\subsection{Tablas y Figuras}
Las tablas y las figuras representan la parte más importante de los resultados de un experimento en particular. Por ello es imprescindible que las mismas sean presentadas con todo el cuidado y atención que sea posible. Es importante indicar que no existe ningún formato preferencial de presentación de tablas y figuras, sin embargo, a fin de lograr una buena percepción de la información suministrada ciertas normas deben ser preservadas independientemente del formato utilizado. Estas normas de presentación no son el resultado de una tendencia o modalidad especifica si no que simplemente, en su conjunto, van dirigidas hacia la presentación clara, completa, lógica y lo mas sencilla posible de los resultados. Existen una innumerable cantidad de recomendaciones a fin de lograr estos objetivos y probablemente la experiencia sea la mejor forma de aprendizaje, en todo caso cualesquiera sea el formato que se escoja, siempre se deben respetar los siguientes lineamientos generales:

\begin{itemize}
\item Expresar claramente las variables que se están evaluando con las unidades respectivas.
\item Respetar los órdenes de dependencia de las variables.
\item Presentar la notación numeral y/o escala más conveniente al rango de las variables.
\item Indicar los errores asociados a los valores de las variables medidas.
\item Respetar el número de cifras significativas y de redondeo.
\item Indicar claramente la naturaleza de las curvas representadas.
\item Evitar el exceso de información.
\item Presentar títulos, leyendas, pies de figura donde se indique claramente la información suministrada.
\item Si se busca realizar comparación directa entre dós curvas o lineas de tendencia, se recomienda ponerlas júntas en un mismo gráfico, diferenciándolas adecuadamente.
\end{itemize}

Ya que pueden ser innumerables los formatos de tablas y figuras que cumplan con estas características, una manera útil de conseguir el formato mas “apropiado” puede obtenerse realizando una revisión de los formatos impuestos en las publicaciones científicas o técnicas de circulación internacional disponibles.

Los títulos de las figuras se colocan luego de las mismas, mientras que los correspondientes a las tablas deben presentarse antes. En ambos casos, deben presentarse en un tamaño de letra o con un interlineado menor al del texto, para lograr su fácil diferenciación del mismo. A continuación dos ejemplos son mostrados: 

\begin{figure}[H]
\centering
 \caption{Aeronave.}
 \includegraphics[width=0.7\textwidth]{aeronave.jpg}
  \label{fig:Dimred}
 \end{figure}

La información de la tabla presentada a continuanción ha sido extraída de \cite{matthews2001aeronautical},

\begin{table}[H]
\centering
\caption{My caption}
\label{my-label}
\begin{tabular}{|c|c|c|c|}
\hline
Liquid       & $kg/m^{3}$ & $lb/ft^{3}$ & Specific gravity \\ \hline
Water        & 1000 & 62,43 & 1                \\ \hline
Sea water    & 1025 & 63,99 & 1,025            \\ \hline
Jet fuel JP1 & 800  & 49,9  & 0,8              \\ \hline
Kerosine     & 820  & 51,2  & 0,82             \\ \hline
Alcohol      & 801  & 50    & 0,801            \\ \hline
Gasoline     & 720  & 44,9  & 0,72             \\ \hline
Oil          & 890  & 55,56 & 0,89             \\ \hline
\end{tabular}
\end{table}

\subsection{Errores}
Todo trabajo experimental esta sujeto a errores. Estos errores se pueden deber a una mala calibración en los instrumentos de medición, a errores de lectura que dependen del experimentador. En todo caso, los errores asociados al trabajo experimental no constituyen un objetivo en si, por lo tanto, no debe concentrarse la discusión de resultados en este aspecto. Esto no quiere decir que no ameriten ser mencionados.

 
El porcentaje de error se debe hallar teniendo en cuenta los datos y resultados. Se relacionan los resultados, se hace una comparación entre el dato teórico y el resultado obtenido a partir de la práctica de laboratorio. Es una análisis objetivo de lo que ocurre experimentalmente y su demostración si es el caso. Se indica de que variables depende el experimento y como afecta su cambio el resultado final, al compararse con un dato teórico.

 
Para el análisis de resultados, se debe calcular el porcentaje de error de la siguiente forma:

\begin{equation}
\% Error = \abs{\frac{(Dato\; te\acute{o}rico)-(Resultado\; experimental)}{(Dato\; te\acute{o}rico)}\times 100}
\end{equation}

\textbf{Nota: El porcentaje de error se halla en valor absoluto.}

El anterior porcentaje, sirve como base para proponer las causas de error que se presentaron durante la práctica de laboratorio y la propagación de errores.

\subsection{Unidades}
Deben ser de preferencia unidades del sistema internacional o bien del sistema métrico. También son aceptables las unidades del sistema ingles, aunque debe tratarse en todo momento de utilizar \textbf{una} en particular durante toda la práctica y el reporte. En cualquier caso hay que ser rigurosos en los siguientes sentidos:

\begin{enumerate}
\item Que las unidades sean correctas, que todas las figuras, gráficas y tablas, especifiquen claramente las unidades en que se has hecho las mediciones o se expresen los resultados.
\item Debe presentarse atención a los cambios de unidades, preguntándose cada vez si tienen sentido los números que se obtienen de cada cambio.
\end{enumerate}

\subsection{Formato de entrega de preinformes e informes}
El informe y el preinforme se entregarán en un mismo documento, el preinforme se entregará antes de laa práctica y el informe completo \textbf{24 horas} después de realizar esta.

Tanto los preinformes como los informes del laboratorio se escribirán usando el compositor de texto \LaTeX\; a través de la herramienta de escritura de textos Overleaf. Será responsabilidad del estudiante aprender el manejo adecuado de este software.

Para hacer esto el docente indicará a los estudiantes el template que deben usar para escribir el documento, los estudiantes lo descargarán y compartirán el documento con su docente para que el docente pueda acceder al documento para revisar los preinformes y los informes.

\end{document}